% ECONOMICS_PROBLEM: {"id": "G23-281", "title": "Shadow banking and system-wide liquidity", "jel_code": "G23", "source_id": "OP-0096"}
% !TeX program = xelatex
\documentclass[11pt,a4paper]{article}
\usepackage[margin=23mm]{geometry}
\usepackage{fontspec}
\setmainfont{Latin Modern Roman}
\usepackage{amsmath,amssymb,mathtools}
\usepackage{xcolor,enumitem,booktabs,longtable,array,xurl,hyperref}
\definecolor{muted}{HTML}{53636C}
\hypersetup{colorlinks=true,urlcolor=blue,linkcolor=blue}
\setlength{\parindent}{0pt}
\setlength{\parskip}{5pt}
\newcommand{\R}{\mathbb R}
\newcommand{\E}{\mathbb E}
\newcommand{\N}{\mathbb N}
\newcommand{\ind}{\mathbf 1}
\newcommand{\Var}{\operatorname{Var}}
\newcommand{\Cov}{\operatorname{Cov}}
\newcommand{\supp}{\operatorname{supp}}
\DeclareMathOperator*{\argmin}{arg\,min}
\DeclareMathOperator*{\argmax}{arg\,max}
\newcommand{\entrymeta}[3]{{\sffamily\small\color{muted}\textbf{\texttt{#1}}\quad|\quad #2\par #3\par}}
\newcommand{\Problemref}[1]{\texttt{#1}}
\newcommand{\JELref}[2]{\texttt{#2} (JEL \texttt{#1})}
\begin{document}
\section*{Shadow banking and system-wide liquidity}
\textbf{Problem ID:} \texttt{G23-281}\quad \textbf{JEL:} \texttt{G23}\par
% BEGIN ECONOMICS STATEMENT
\label{problem:OP-0096}
\label{atlas:prob:F6}

\noindent\textbf{Atlas ID:} F6; \textbf{original number:} 96; \textbf{field:} Finance. \textbf{Classification:} Editorial research formulation (theoretical/empirical); not a certified unsolved theorem.

\subsection*{Definitions and assumptions}
Let financial entities $i=1,\ldots,n$ include banks, dealers and funds. Model $\theta$ specifies assets, maturity-dated liabilities, collateral, repo haircuts, derivatives and legally enforceable netting sets. Balance-sheet identity $A_i=L_i+E_i$ separates asset value $A_i$, liabilities $L_i$ and equity $E_i$. A liquidity shock $Z$ changes redemption and margin needs. Entities choose feasible lending, funding and asset sales subject to solvency, collateral and payment constraints; market prices and contractual payments jointly satisfy a selected clearing equilibrium. Regulation $r$ changes constraints and the endogenous pre-shock allocation. Let $\ell(Z,r)$ be finite total real loss to ultimate investors and borrowers.

\subsection*{Formal research question}
Identify the regulation-induced redistribution of short-term funding and collateral exposures and its effect on
\[
 \mathcal R(r)=\mathbb E_\theta[\ell(Z,r)]
 \quad\text{and}\quad
 T_q(r)=\inf\{x:\mathbb P_\theta(\ell(Z,r)\le x)\ge q\},
\]
where $q\in(0,1)$ is a fixed loss-quantile level. Characterize conditions under which a tightening of bank regulation lowers bank fragility while increasing system-wide losses through nonbank substitution. Determine the minimum legal-entity, maturity and collateral information that identifies exposures and supports meaningful bounds when parts of the network are unobserved.

\subsection*{Interpretation and scope}
Leverage ratios and gross notional derivative values are not directly comparable measures of liquidity risk. Collateral reuse, netting and common underlying assets can make summing bilateral exposures double count claims or misstate cash needs. Haircut increases and asset sales feed back through market prices; a static stress test with fixed holdings omits endogenous migration before regulation and adjustment after shocks. A loss quantile requires a stated shock law and may be discontinuous; do not infer it from a single severe scenario. Welfare analysis must distinguish transfers between intermediaries from real liquidation, bankruptcy and credit-supply costs. Taxpayer guarantees and emergency lending belong in the policy environment.

\subsection*{Source audit and status}
[S1] concerns leverage cycles, [S2] repo runs, and [S3] documents shadow-banking funding chains; [S3] first appeared in 2010 and was revised in February 2012. Atlas link [49], [S4], is an institutional homepage supporting monitoring context rather than one precisely located result. This is an editorial system-wide measurement and causal-policy agenda, not a certified open clearing theorem.

\subsection*{References}

\begin{enumerate}[label={[S\arabic*]},leftmargin=*]

\item Tobias Adrian and Hyun Song Shin (2010). \emph{Liquidity and Leverage}. Journal of Financial Intermediation 19(3), 418–437. \url{https://www.newyorkfed.org/research/staff_reports/sr328.html}.
\textit{Locator and role:} Publisher or author abstract / bibliographic record; Procyclical leverage and intermediary balance sheets. Provides background or a benchmark result; does not establish that the editorial question remains unresolved in 2026.

\item Gary Gorton and Andrew Metrick (2012). \emph{Securitized Banking and the Run on Repo}. Journal of Financial Economics 104(3), 425–451. \url{https://www.sciencedirect.com/science/article/pii/S0304405X1100081X}.
\textit{Locator and role:} Publisher or author abstract / bibliographic record; Repo haircuts and liquidity stress during the financial crisis. Provides background or a benchmark result; does not establish that the editorial question remains unresolved in 2026.

\item Zoltan Pozsar, Tobias Adrian, Adam Ashcraft, and Hayley Boesky (2010). \emph{Shadow Banking}. Federal Reserve Bank of New York Staff Report 458; revised February 2012. \url{https://www.newyorkfed.org/research/staff_reports/sr458.html}.
\textit{Locator and role:} Publisher or author abstract / bibliographic record; Taxonomy and funding-flow maps of shadow banking; original 2010 version is available. The webpage also exposes a revised February 2012 version. Version dates are not interchangeable.

\item European Systemic Risk Board (2026). \emph{European Systemic Risk Board}. Institutional homepage, accessed 8 October 2026. \url{https://www.esrb.europa.eu/}.
\textit{Locator and role:} Publisher or author abstract / bibliographic record; Homepage identifies current system-wide liquidity and bank/nonbank linkage research. An institutional homepage provides context rather than a precisely specified theorem or empirical result; archived report citation is preferable for individual findings.

\end{enumerate}

\subsection*{Common conventions: Source mathematical conventions}

Unless an entry states otherwise, agent and firm sets are finite. State and action spaces are measurable, strategies and outcome maps are measurable, probability laws are specified, and expectations exist for the targets under discussion. Infinite-horizon discount factors lie in $(0,1)$ and discounted objectives are finite. Norms, tolerances and complexity input sizes must be specified for computational comparisons. Constants or primitives that appear locally are inputs, not empirical facts asserted by this catalogue.

\textbf{Identification.} A statistical model $m\in\mathcal M$ implies an observable law $P_m$ and target $\theta(m)$. At law $P$, its identified set is
\[
\Theta(P;\mathcal M)=\{\theta(m):m\in\mathcal M,\ P_m=P\}.
\]
Point identification holds when this set is a singleton, or equivalently when $P_{m_1}=P_{m_2}$ implies $\theta(m_1)=\theta(m_2)$ for all compatible models. Sharp bounds mean bounds attained, or approached when the boundary is not attained, by compatible models. Writing this set is a definition, not a solution: the substantive task is to establish informative restrictions and derive or estimate the set. An unrestricted-confounding model may give only trivial bounds.

\textbf{Inference.} Identification, estimability and valid uncertainty quantification are separate questions. Fix $\alpha\in(0,1)$. A confidence procedure $C_n$ over a declared law class $\mathcal P$ has asymptotic uniform coverage at level $1-\alpha$ when
\[
\liminf_{n\to\infty}\inf_{P\in\mathcal P}\Pr_P\{\theta(P)\in C_n\}\geq1-\alpha.
\]
For set-identified targets the coverage event must be defined separately, for example coverage of the full identified set. If an entry uses identification-indexed models instead of uniquely indexed laws, the same coverage convention is applied over those models. Pointwise limits do not establish uniform validity, and asymptotic validity does not guarantee finite-sample precision. Sample sizes, dependence structures and weak-identification sequences are part of each application.

\textbf{Counterfactuals.} $Y_i(a)$ is a potential outcome under a specified intervention, including its timing, implementation and versions. It is not $Y_i$ conditional on voluntarily choosing $a$. A full intervention vector permits interference; shorthand scalar treatment notation does not silently impose no interference. Assignment, selection, attrition, anticipation and measurement must be specified. A claim about an instrument's local effect is not automatically a population policy effect.

\textbf{Economic equilibrium.} A counterfactual model includes best responses, constraints, feasibility, market clearing, state transitions and expectations consistent with the stated information. When several equilibria exist, either a declared selector is an input or the target is set-valued. Comparisons across policy regimes must use the stated selection convention. Reduced-form relationships are not presumed invariant after an equilibrium-changing intervention.

\textbf{Optimization and welfare.} Optimization is over the stated feasible set. If attainment is not established, $\sup$ or $\inf$ and $\varepsilon$-optimal choices replace an asserted maximizer or minimizer. For example, with value $V=\sup_{a\in\mathcal A}W(a)<\infty$, an $\varepsilon$-optimal set is $\{a\in\mathcal A:W(a)\geq V-\varepsilon\}$ for $\varepsilon>0$. Benefits, costs, risk and health or environmental outcomes can be aggregated only using declared units and conversion weights. Interpersonal utility weights, the treatment of fiscal transfers and foreign profits, discounting, and the behavioral welfare criterion are explicit inputs. A comparative-static derivative additionally requires existence and appropriate differentiability of the selected counterfactual.

\textbf{Sources.} Local labels [S1], [S2], and so forth restart in every question. Locators identify the relevant abstract, model, section or result at the level actually checked; a bibliographic or abstract-level verification is not represented as inspection of an entire proof. Working papers are distinguished from later publications and changing versions. Source summaries are paraphrased. All URL checks and editorial formulations refer to the audit date stated above.
% END ECONOMICS STATEMENT
\end{document}
